\documentclass{article}
\usepackage{amsmath, amssymb, amsthm}
\usepackage{hyperref}
\usepackage{geometry}
\geometry{margin=1in}

\title{Erd\H{o}s Problem \#422}
\date{Last updated: 2025-08-31}
\author{Source: erdosproblems.com}

\begin{document}
\maketitle

\noindent\textbf{Status:} OPEN \quad \textbf{No prize}

\noindent\textbf{Tags:} number theory

\medskip
\noindent\textbf{Problem Statement:}

Let $f(1)=f(2)=1$ and for $n>2$\[f(n) = f(n-f(n-1))+f(n-f(n-2)).\]Does $f(n)$ miss infinitely many integers? What is its behaviour?

\bigskip
\noindent\textbf{Remarks:}

Asked by Hofstadter. The sequence begins $1,1,2,3,3,4,\ldots$ and is A005185 in the OEIS. It is not even known whether $f(n)$ is well-defined for all $n$.

\bigskip
\noindent\small{Source: \url{https://www.erdosproblems.com/422}}
\end{document}
